% Generated by paperkit. Do not edit.
\newcommand{\PaperTitle}{When Should a Shopping Agent Stop Searching?}
\newcommand{\PaperSubtitle}{A measured decision rule for adaptive stopping under hard budgets}
\newcommand{\PaperAuthors}{Ahnaf Prio\\{\small Independent Researcher}}
\newcommand{\PaperAbstract}{A shopping agent has to decide after each merchant: attempt to buy the current catalog offer or pay to keep searching without reserving it. A fixed threshold is easy to deploy, but it is unclear when a budget-aware decision can save enough to justify the time, calls, and spend of another search. We answer that question with live catalog-search data from independent merchants. Across two dated snapshots, listings disappear and surviving prices change within a month, while repeated scans minutes apart are byte-identical. We replay eight stopping policies on frozen held-out panels from those snapshots. An adaptive reservation-price rule lowers failure-penalized purchase loss relative to the strongest fixed rule tuned on calibration data; the paired bootstrap interval remains below zero, and every replay respects its hard budget. The improvement is statistically detectable but small because most matched products have only two merchants and little price variation. That limitation is the main finding, not a footnote: a calibrated sweep shows that adaptive stopping helps when merchants differ enough in price and the budget leaves room to use another observation. It is indistinguishable from a tuned fixed rule when prices are nearly identical and includes counterexamples where a fixed rule wins. Retained offers and revalidation are outside the current model. We release the optimizer, replay harness, and collected panels.}
\newcommand{\PaperKeywords}{autonomous shopping, online search, optimal stopping, resource constraints, ephemeral offers}
\newcommand{\PaperArxivCategories}{cs.AI}
